\documentclass[7pt]{article}
%\usetheme{Warsaw}

%\usepackage{beamerinnerthemecircles, beamerouterthemeshadow}
%\usepackage{scrextend}

%\changefontsizes{7pt}


% Standard packages
\usepackage{graphicx}
\usepackage{graphics}

\usepackage{color}
\usepackage[OT4]{fontenc}
\usepackage[utf8]{inputenc}
%\usepackage[latin1]{inputenc}
\usepackage{times}\usepackage{amsmath}
\usepackage[T1]{fontenc}

%\mode<presentation> {
%    \usetheme{Singapore} %Bergen, Berkely, Berlin, Boadilla, CambridgeUS, Darmstadt,
%                          %Frankfurt, Goettingen, Singapore, Warsaw
%    \usecolortheme{default} %beetle, seahorse, wolverine, dolphin, beaver
%}


\title{Forced oscillations for an electron in a 2D square infinite quantum well}

\definecolor{whi}{rgb}{0.95,0.95,0.95}
% Setup TikZ

%\usepackage{tikz}
%\usetikzlibrary{arrows}
%\tikzstyle{block}=[draw opacity=0.7,line width=1.4cm]


% Author, Title, etc.




% The main document

\begin{document}
\maketitle




\section{the problem}

We return to the problem of an electron in a 2D infinite quantum well (three assignments back).
We will try to simulate the transitions driven by an external electric field. 
For that purpose we introduce a time-dependent Hamiltonian $H(t)=H_0+eFx \sin\omega t$,
where $H_0$ is the stationary Hamiltonian from the previous assignement and the second term introduced
the electric field potential with the amplitude eF, oriented along the  $x$ direction and changing
with the freqency of $\omega$. 
The dynamics is governed by the Schroedinger equation
\begin{equation} i\hbar \frac{\partial \Psi}{\partial t}=H(t)\Psi\end{equation}
For the initial moment in time we set the ground-state -- that you calculated three assignments ago.

\subsection{potential matrix}
Now, we need to include the electric potential to the Hamiltonian. We will need a global matrix 
of a {\it uniform} electric field $eFx$. Each element has its own local matrix of this potential. The local
matrix of this potential is different in each element -- unlike the kinetic energy operator. 

We calculate the local potential mattrices with 
$$v^k_{ji}=\frac{a^2}{4} \frac{m\omega^2}{2}\int_{\Omega_k} eF x(\xi_1) g_j(\xi_1,\xi_2) g_i(\xi_1,\xi_2)d\xi_1 d\xi_2.$$
We use the Gauss quadrature as previously, only now wee need to evaluate $x$ for the Gauss points
and they will be different in each element.
$$v^k_{ji}=\frac{a^2}{4}\frac{m\omega^2}{2}\sum_{l=1}^3\sum_{n=1}^3  w_l w_n eF x(p_l) g_j(p_l,p_n) g_i(p_l,p_n)$$
The evaluation of $x$ proceeds as previously
\begin{eqnarray}
x(\xi_1,\xi_2)&=&\sum_{i=1}^4 x_{nlg(k,i)}g_i(\xi_1,\xi_2).\label{wzor} 
\end{eqnarray}
After evaluation of the local matrix elements we produce the global matrix $V$ -- in the manner that we used
for our previous FEM projects.

Our wave function is still expanded in the basis of shape functions but the coefficients $c_i$ now depend on time
   \begin{equation} \Psi=\sum_{i=1}^N c_i(t) h_i(x,y).\end{equation} 
The Crank-Nicolson scheme for our problem produces a linear system of equations for $c_i(t+\Delta t)$ that needs to be solved for each time step.
\begin{equation}
\Psi(t+\Delta t)=\Psi(t)+\frac{\Delta t}{2\hbar i} \left(H(t+\Delta t)\Psi(t+\Delta t)+H(t)\Psi(x,t)\right).
\end{equation}
Upon substitution of the expansion in the basis we get 
{
\begin{equation}
\left[{\bf S}-\frac{\Delta t}{2\hbar i} {\bf H}(t+\Delta t)\right]{\bf c}(t+\Delta t)=\left[{\bf S}+\frac{\Delta t}{2\hbar i}{\bf H}(t)\right]{\bf c}(t),
\end{equation}
}
where ${\bf H}(t)={\bf H}_0+{\bf V} \sin(\omega t)$.

Lets set $\omega=\Delta E/\hbar$ where $\Delta E=E_1-E_0$ and $eF=1/L$, $\Delta t=0.25$.
We need to monitor (1) the conservation of the norm of the wave function
$N(t)={\bf c(t)}^T {\bf O} {\bf c(t)}$ where ${\bf O}$ is the overlap matrix
and (2) 
the projections of the wave function on the $H_0$ eigenstates
$p_i(t)=\langle \Psi_i | \Psi(t)\rangle={\bf c}_i^T {\bf O} {\bf c(t)}$.
The probability that the electron occupies eigenstate $i$ in the moment in time $t$ is $|p_i|^2$.

\section{Deliverables}
\begin{itemize}
\item Plot $|p_i(t)|^2$ for the lowest-energy states. What is the probability of the electron transition probability to one of the two degenerate first
excited states ?
\item Replace $\omega$ but half of the nominal resonant value. What is the result now? 
\item get back to the resonant $\omega$. What is the leakage of the wave function to the higher energy states
than the ground-state and the higher excited states and how the leakage depends on the amplitude of the field.
\end{itemize}
\end{document}




